\section{ISR 汇编 Stub：\texttt{isr\_stubs.asm}}

这是 IDT 实现中最精密的部分——必须保证汇编 push 的顺序与 C 结构体\textbf{完全一致}。

\subsection{问题：栈布局不统一}

当异常发生时，CPU 自动压栈：

\begin{figure}[H]
  \centering
  % figures/ch04/fig03-cpu.tex — auto-extracted from chapter source
\begin{tikzpicture}[node distance=0cm]
  % CPU 自动压栈（从高到低）
  \node[memcell, minimum width=5cm, fill=yellow!15] (efl) {EFLAGS};
  \node[memcell, minimum width=5cm, fill=yellow!15, below=0cm of efl] (cs) {CS};
  \node[memcell, minimum width=5cm, fill=yellow!15, below=0cm of cs] (eip) {EIP};
  \node[memcell, minimum width=5cm, fill=red!15, below=0cm of eip] (err) {err\_code（仅部分异常）};

  \node[right=0.5cm of efl, font=\scriptsize\itshape] {CPU 自动压入};
  \node[right=0.5cm of err, font=\scriptsize\itshape, text=red!60!black] {只有向量 8,10-14,17 有};

  \node[memaddr, left=0.3cm of err] {ESP →};

  % 高地址标注
  \node[above=0.3cm of efl, font=\scriptsize] {高地址 ↑};
  \node[below=0.3cm of err, font=\scriptsize] {低地址 ↓};
\end{tikzpicture}

  \caption{异常入口时 CPU 自动压栈的内容}
\end{figure}

问题：有些异常（如 \#PF）CPU 会额外压入一个 error code，有些（如 \#DE）不会。
如果不处理这个差异，C handler 无法用统一的结构体解析栈。

\subsection{解决方案：ISR\_NOERR 和 ISR\_ERR 宏}

\codefile{src/kernel/arch/isr\_stubs.asm}
\begin{lstlisting}[style=nasmstyle]
; 无错误码的异常：手动 push 0 对齐
%macro ISR_NOERR 1
isr%1:
    push dword 0          ; 伪 err_code = 0
    push dword %1         ; int_no
    jmp isr_common_stub
%endmacro

; 有错误码的异常：CPU 已经 push 了 err_code
%macro ISR_ERR 1
isr%1:
    push dword %1         ; int_no（err_code 已在栈上）
    jmp isr_common_stub
%endmacro
\end{lstlisting}

效果：无论哪种异常，进入 \texttt{isr\_common\_stub} 时栈顶布局总是统一的：

\begin{figure}[H]
  \centering
  % figures/ch04/fig04-isr_common_stub.tex — auto-extracted from chapter source
\begin{tikzpicture}[node distance=0cm]
  \node[memcell, minimum width=5cm, fill=yellow!15] (efl) {EFLAGS};
  \node[memcell, minimum width=5cm, fill=yellow!15, below=0cm of efl] (cs) {CS};
  \node[memcell, minimum width=5cm, fill=yellow!15, below=0cm of cs] (eip) {EIP};
  \node[memcell, minimum width=5cm, fill=orange!15, below=0cm of eip] (err) {err\_code};
  \node[memcell, minimum width=5cm, fill=orange!15, below=0cm of err] (ino) {int\_no};

  \node[right=0.5cm of efl, font=\scriptsize\itshape] {CPU 自动压入};
  \node[right=0.5cm of err, font=\scriptsize\itshape] {stub 确保存在};
  \node[right=0.5cm of ino, font=\scriptsize\itshape] {stub push};

  \node[memaddr, left=0.3cm of ino] {ESP →};
\end{tikzpicture}

  \caption{进入 isr\_common\_stub 时的统一栈布局}
\end{figure}

\subsection{\texttt{isr\_common\_stub}：逐指令详解}

\codefile{src/kernel/arch/isr\_stubs.asm}
\begin{lstlisting}[style=nasmstyle, numbers=left]
isr_common_stub:
    pushad              ; (1) 保存 8 个通用寄存器

    push ds             ; (2) 保存段寄存器
    push es             ; (3)
    push fs             ; (4)
    push gs             ; (5)

    mov ax, 0x10        ; (6) 切换到内核数据段
    mov ds, ax          ; (7)
    mov es, ax          ; (8)
    mov fs, ax          ; (9)
    mov gs, ax          ; (10)

    push esp            ; (11) 参数：指向 regs_t 的指针
    call isr_handler_c  ; (12) 调用 C handler
    add esp, 4          ; (13) 清理参数

    pop gs              ; (14) 恢复段寄存器
    pop fs              ; (15)
    pop es              ; (16)
    pop ds              ; (17)

    popad               ; (18) 恢复通用寄存器

    add esp, 8          ; (19) 跳过 int_no 和 err_code

    iret                ; (20) 从中断返回
\end{lstlisting}

\paragraph{第 2 行：\texttt{pushad}}

\texttt{pushad} 一次性压入 8 个通用寄存器，顺序是：
EAX → ECX → EDX → EBX → ESP → EBP → ESI → EDI。

\begin{warnbox}
\texttt{pushad} 压入的 ESP 是\textbf{执行 pushad 之前}的值（即"当时的快照"），
不是当前栈指针。这意味着 \texttt{regs\_t} 中的 \texttt{esp} 字段只是参考值，
不能用来恢复栈。\texttt{popad} 恢复时会跳过 ESP 字段。
\end{warnbox}

\paragraph{第 3--6 行：保存段寄存器}

\texttt{push ds/es/fs/gs} 的顺序决定了 \texttt{regs\_t} 中对应字段的位置。
\textbf{顺序必须严格匹配}——下面会看到 C 结构体的字段排列。

\paragraph{第 7--11 行：切换到内核数据段}

异常可能发生在段寄存器指向非内核段的时候（将来用户态会遇到）。
在调用 C handler 前，先确保 DS/ES 指向内核数据段 \hex{10}。

\paragraph{第 12--14 行：调用 C handler}

\texttt{push esp} 把当前栈指针作为参数传给 \texttt{isr\_handler\_c}。
此时 ESP 恰好指向刚才 push 的所有数据的起始位置——也就是 \texttt{regs\_t*}。

\paragraph{第 20 行：\texttt{add esp, 8}}

跳过栈上的 \texttt{int\_no} 和 \texttt{err\_code}（各 4 字节），
让 ESP 指向 CPU 自动压入的 EIP/CS/EFLAGS。

\paragraph{第 21 行：\texttt{iret}}

\texttt{iret}（Interrupt Return）从栈上弹出 EIP、CS、EFLAGS，恢复到中断前的执行状态。

\begin{cpustate}
\texttt{iret} 执行时 CPU 弹出三个值：

\begin{enumerate}
  \item 弹出 EIP → 恢复程序计数器
  \item 弹出 CS → 恢复代码段（如果 CS.RPL 变了，触发特权级切换）
  \item 弹出 EFLAGS → 恢复标志位（包括 IF，如果原来中断是开的就恢复开启）
\end{enumerate}

如果发生了特权级切换（如从 Ring 0 回到 Ring 3），还会额外弹出 ESP 和 SS。
\end{cpustate}

% ============================================================================

